% Gautama Documentation - LaTeX Book Template
% Professional book layout for technical documentation

\documentclass[11pt,a4paper,twoside,openright]{book}

% ============================================================================
% Package Imports
% ============================================================================

% Core packages
\usepackage[utf8]{inputenc}
\usepackage[T1]{fontenc}
\usepackage[english]{babel}

% Page layout
\usepackage[
    top=2.5cm,
    bottom=2.5cm,
    left=3cm,
    right=2.5cm,
    headheight=15pt
]{geometry}

% Fonts
\usepackage{lmodern}           % Latin Modern font
\usepackage{inconsolata}       % Better monospace font for code
\usepackage{microtype}         % Better typography

% Colors
\usepackage[dvipsnames]{xcolor}
\definecolor{linkcolor}{RGB}{0,102,204}
\definecolor{codebackground}{RGB}{245,245,245}
\definecolor{codeborder}{RGB}{200,200,200}

% Graphics and figures
\usepackage{graphicx}
\usepackage{float}
\usepackage{caption}
\usepackage{subcaption}

% Tables
\usepackage{booktabs}
\usepackage{longtable}
\usepackage{array}

% Lists
\usepackage{enumitem}

% Links and references
\usepackage[
    colorlinks=true,
    linkcolor=linkcolor,
    urlcolor=linkcolor,
    citecolor=linkcolor,
    bookmarks=true,
    bookmarksnumbered=true,
    pdfstartview=FitH
]{hyperref}

% Code listings
\usepackage{listings}
\usepackage{fancyvrb}

% Headers and footers
\usepackage{fancyhdr}

% Title page
\usepackage{titling}

% Table of contents
\usepackage{tocloft}

% Index
\usepackage{makeidx}
\makeindex

% ============================================================================
% Code Listing Configuration
% ============================================================================

\lstdefinestyle{codestyle}{
    backgroundcolor=\color{codebackground},
    commentstyle=\color{gray},
    keywordstyle=\color{blue}\bfseries,
    numberstyle=\tiny\color{gray},
    stringstyle=\color{orange},
    basicstyle=\ttfamily\small,
    breakatwhitespace=false,
    breaklines=true,
    captionpos=b,
    keepspaces=true,
    numbers=left,
    numbersep=5pt,
    showspaces=false,
    showstringspaces=false,
    showtabs=false,
    tabsize=2,
    frame=single,
    rulecolor=\color{codeborder},
    xleftmargin=10pt,
    xrightmargin=10pt
}

\lstset{style=codestyle}

% Language definitions
\lstdefinelanguage{nix}{
    keywords={let, in, with, inherit, import, if, then, else, assert, rec, or},
    sensitive=true,
    morecomment=[l]{\#},
    morestring=[b]",
}

% ============================================================================
% Page Style Configuration
% ============================================================================

\pagestyle{fancy}
\fancyhf{}

% Headers
\fancyhead[LE]{\nouppercase{\leftmark}}
\fancyhead[RO]{\nouppercase{\rightmark}}

% Footers
\fancyfoot[LE,RO]{\thepage}
\fancyfoot[RE,LO]{\small Gautama Documentation}

% Plain style for chapter pages
\fancypagestyle{plain}{
    \fancyhf{}
    \fancyfoot[LE,RO]{\thepage}
    \renewcommand{\headrulewidth}{0pt}
}

% ============================================================================
% Title Page Configuration
% ============================================================================

\pretitle{%
  \begin{center}
  \LARGE
  \includegraphics[width=6cm]{logo.png}\\[\bigskipamount]
}
\posttitle{\end{center}}

\title{\Huge\bfseries Gautama\\System Documentation}
\author{%
    System Administrator\\[0.5cm]
    \large NixOS Production Infrastructure
}
\date{\today}

% ============================================================================
% Custom Commands
% ============================================================================

% Inline code
\newcommand{\code}[1]{\texttt{#1}}

% File path
\newcommand{\filepath}[1]{\texttt{\textbf{#1}}}

% Command
\newcommand{\cmd}[1]{\texttt{\$ #1}}

% Environment variable
\newcommand{\env}[1]{\texttt{\textbf{#1}}}

% URL styling
\newcommand{\urllink}[2]{\href{#1}{\color{linkcolor}\underline{#2}}}

% Warning box
\usepackage{tcolorbox}
\newtcolorbox{warningbox}{
    colback=red!5!white,
    colframe=red!75!black,
    title=Warning,
    fonttitle=\bfseries
}

% Info box
\newtcolorbox{infobox}{
    colback=blue!5!white,
    colframe=blue!75!black,
    title=Information,
    fonttitle=\bfseries
}

% Tip box
\newtcolorbox{tipbox}{
    colback=green!5!white,
    colframe=green!75!black,
    title=Tip,
    fonttitle=\bfseries
}

% ============================================================================
% Chapter and Section Styling
% ============================================================================

\usepackage{titlesec}

% Chapter style
\titleformat{\chapter}[display]
    {\normalfont\huge\bfseries\color{MidnightBlue}}
    {\chaptertitlename\ \thechapter}
    {20pt}
    {\Huge}

% Section style
\titleformat{\section}
    {\normalfont\Large\bfseries\color{MidnightBlue}}
    {\thesection}
    {1em}
    {}

% ============================================================================
% Table of Contents Styling
% ============================================================================

\renewcommand{\cftchapfont}{\bfseries\color{MidnightBlue}}
\renewcommand{\cftsecfont}{\color{black}}

% ============================================================================
% Document Begin
% ============================================================================

\begin{document}

% Front matter
\frontmatter

% Title page
\maketitle

% Copyright page
\newpage
\thispagestyle{empty}
\vspace*{\fill}
\begin{center}
    \textbf{Gautama System Documentation}\\[0.5cm]
    Copyright \copyright\ \the\year\ System Administrator\\[0.5cm]
    All rights reserved.\\[1cm]

    \textit{This documentation covers the complete Gautama NixOS infrastructure,\\
    including all services, deployment procedures, and operational guidelines.}\\[1cm]

    \textbf{Generated:} \today\\[0.5cm]
    \textbf{Version:} 1.0\\[0.5cm]
    \textbf{Format:} LaTeX PDF
\end{center}
\vspace*{\fill}

% Table of Contents
\tableofcontents

% List of Figures
\listoffigures

% List of Tables
\listoftables

% Preface
\chapter*{Preface}
\addcontentsline{toc}{chapter}{Preface}

This documentation provides comprehensive coverage of the Gautama NixOS infrastructure,
a production-grade self-hosted system running on Apple Silicon hardware.

The system encompasses over 60 services, including databases, monitoring, networking,
containers, and web applications, all managed declaratively through NixOS.

\textbf{Key Features:}
\begin{itemize}
    \item Declarative configuration with NixOS flakes
    \item Reproducible builds and deployments
    \item Comprehensive monitoring with Prometheus and Grafana
    \item Secure secret management with SOPS
    \item Container orchestration with Podman Quadlet
    \item Private certificate authority with step-ca
    \item ZFS storage with automatic snapshots
    \item Cloud backups to Backblaze B2
\end{itemize}

\textbf{Audience:}

This documentation is intended for system administrators, DevOps engineers, and
technical users responsible for maintaining, deploying, or extending the Gautama
infrastructure.

\textbf{Structure:}

The documentation is organized by topic, with each chapter covering a specific
aspect of the system. Cross-references are provided throughout for related topics.

% Main matter
\mainmatter

%%% CONTENT WILL BE INSERTED HERE BY BUILD SCRIPT %%%

% Back matter
\backmatter

% Index
\printindex

% Colophon
\chapter*{Colophon}
\addcontentsline{toc}{chapter}{Colophon}

This book was typeset using \LaTeX, a document preparation system based on \TeX.

\textbf{Tools Used:}
\begin{itemize}
    \item \textbf{Typesetting:} \LaTeX\ with pdflatex
    \item \textbf{Conversion:} Pandoc for Markdown to \LaTeX
    \item \textbf{Build System:} Nix with automatic rebuilds
    \item \textbf{Fonts:} Latin Modern (text), Inconsolata (code)
    \item \textbf{Graphics:} Inkscape, D2 diagrams
\end{itemize}

\textbf{Source:}

The source files for this documentation are maintained in Markdown format in the
\code{/etc/nixos/docs/md/} directory and automatically converted to PDF during
system rebuilds.

\vspace{1cm}

\textit{Built with} \LaTeX\ \textit{on NixOS}

\end{document}
